\chapter{An unfamiliar reader walkthrough}\label{reader-walkthrough}

An unfamiliar-reader session is planned. The following protocol is prepared; no human observations have been recorded. The companion includes an empty recording form at \path{walkthrough/results-template.json}.

Status: \textbf{prepared, not conducted}. No participant results are included. An agent running commands is an automated reproduction check, not an unfamiliar human reader.

\section{Facilitation}

Use someone who has not authored or rehearsed the example. Give them the exact book edition and companion commit, a clean checkout and the task prompts below. Let them follow the book before offering a hint. Record the page, command, observed result and any assistance. A successful assisted step is different from an independent completion. Leave the original confusing passage identifiable when recording its proposed repair.

Start with the Node/npm versions recorded by the edition. Record terminal, operating system, viewport and any assistive tool. \texttt{results-template.json} is an empty recording form, not a verification receipt. Keep \texttt{not\_run} until a real session occurs; do not populate answers with an agent's predictions.

\section{Session A — Build from a clean checkout}

\begin{enumerate}
\item Find the edition's source and reproduce installation without facilitator help.
\item Run stage 01. Explain which values identify a job and which describe its label.
\item Read the first diff and run stage 02. Identify the failed job without color and describe what pipe mode preserves. Notice that stage 02's interactive specimen is a printed fixture, not a live session.
\item Run stage 03. Start compile, open cancellation, attempt movement, dismiss the dialog, then cancel it. Explain the state after each action.
\item Read the asynchronous step. Run stage 04, open compile logs and select test. Explain why a late response must not replace current detail. Use the test or exercise if wall-clock timing does not make the race visible.
\item Run stage 05 and resize to 24×8. Find the current status and available actions. Inspect a failure and return to another job. Explain what information scrolling preserves and what is no longer simultaneously visible.
\item Attempt one broken exercise before using its answer. Use the hint, then the diagnosis if necessary; record which level was needed.
\item Find the license and the exact companion commit used by the book. Identify one claim the automated tests do not establish.
\end{enumerate}

Completion criteria: installation is reproducible; every stage's declared outcome can be observed; the reader can distinguish state, effects and presentation; all required corrections are attached to a reproducible task and passage. Record failures and assistance instead of turning the criteria into assumed passes.

\section{Session B — Keyboard-only interactive task}

Run the final stage in a supported terminal without using a pointer. Ask the participant to start a job, inspect logs, select another job, filter failures, open and dismiss cancellation, and quit. Include a narrow viewport. Record focus confusion, discoverability, hidden controls, misleading status and final terminal state. The completion is a human task observation, not merely a key-routing test.

\section{Session C — Reading with an assistive tool}

Use the participant's actual tool and record its name and version. Run the accessible fixture and the color-free static fixture separately. Ask which job failed, why it failed, what is selected, and which actions the text describes. Observe reading order and ambiguity. Do not substitute the expected answer for what the participant actually understood.

The accessible CLI fixture prints once and exits. It does not implement an interactive screen-reader session. Testing a live screen-reader workflow requires an implemented interaction contract and its own task session. Record that item as \texttt{not\_implemented}, not as a failed or passed snapshot-reading observation.

\section{Turn observations into repairs}

For every problem retain: edition, source commit, task, page, exact command, expected observation, actual observation, assistance, reproduction and proposed change. Distinguish a code defect, an explanation gap, an environment mismatch and an unimplemented capability. After a repair, repeat the affected task; use another unfamiliar participant before calling its discoverability independently verified.
